\documentclass[11pt,a4paper]{article}

\usepackage[T1]{fontenc}
\usepackage[utf8]{inputenc}
\usepackage[shorthands=off,turkish]{babel}
\usepackage[a4paper,margin=2.4cm]{geometry}
\usepackage{lmodern}
\usepackage{booktabs}
\usepackage{longtable}
\usepackage{array}
\usepackage{enumitem}
\usepackage{xcolor}
\usepackage{fancyhdr}
\usepackage[hidelinks]{hyperref}

% --- Renk tanimlari ---
\definecolor{kritik}{HTML}{B91C3C}
\definecolor{onemli}{HTML}{C07807}
\definecolor{teslim}{HTML}{1D6B8E}
\definecolor{tamam}{HTML}{04704D}
\definecolor{acikgri}{HTML}{F3F6F5}

\newcommand{\kritikrozet}{\textcolor{kritik}{\textbf{KRİTİK}}}
\newcommand{\onemlirozet}{\textcolor{onemli}{\textbf{ÖNEMLİ}}}
\newcommand{\teslimrozet}{\textcolor{teslim}{\textbf{TESLİM}}}

% Basit vurgu kutusu (ek pakete ihtiyaç duymaz)
\newenvironment{vurgukutu}[1]
  {\par\vspace{6pt}\noindent
   \begin{tabular}{|@{\hspace{8pt}}p{0.94\textwidth}@{\hspace{8pt}}|}
   \hline\rule{0pt}{16pt}\textcolor{kritik}{\textbf{#1}}\\[4pt]}
  {\\[6pt]\hline
   \end{tabular}\par\vspace{6pt}}

\pagestyle{fancy}
\fancyhf{}
\lhead{\small Katılım Bankacılığı NLP Çözümü}
\rhead{\small TEKNOFEST 2026 --- 2. Senaryo}
\cfoot{\thepage}

\title{\vspace{-1.5cm}\textbf{Eksiklik Analizi ve Yol Haritası}\\[0.3cm]
\large TEKNOFEST 2026 Yapay Zekâ Dil Ajanları Yarışması --- 2. Senaryo\\
\normalsize Katılım Bankacılığı Kampanya Metinlerinden Bilgi Çıkarımı}
\author{}
\date{\today}

\begin{document}
\maketitle
\thispagestyle{fancy}

\vspace{-1cm}

\section{Yönetici Özeti}

Mevcut çözüm, katılım bankacılığı kampanya metinlerinden yapılandırılmış JSON çıkaran
çalışır durumda bir sistemdir. Terim eşleme (faiz $\rightarrow$ kâr payı), sayısal
normalizasyon, alan bazlı güven skoru ve \textbf{her alan için kanıt alıntısı} üretimi
gibi güçlü yanları vardır. Kanıt alıntısı üretimi şartnamede istenmeyen fakat jüri
karşısında ayırt edici olacak bir yetenektir.

Buna karşılık şartnamenin üç temel bileşeninde ciddi boşluk bulunmaktadır:
\textbf{on-premise çalışabilirlik}, \textbf{chatbot} ve \textbf{gerçek veri seti}.
Bu üç eksik, toplam puanın yaklaşık yarısını doğrudan etkilemektedir.

\begin{vurgukutu}{En acil bulgu}
Çözüm her istekte bulut tabanlı NVIDIA NIM servisine gitmektedir. Şartname Madde 5.9
kurum içi çalışabilirliği, müşteri verilerinin kurum dışına çıkmamasını ve dış
servislere bağımlı olmamayı açıkça şart koşmaktadır. Ayrıca Madde 8, projenin bağımlı
olduğu ücretli yazılım ve üçüncü taraf hizmet kullanımını yasaklamaktadır. Bu tek
başına \%20 ağırlıklı ``On-Prem Uygulanabilirlik'' kriterinin neredeyse tamamını
riske atmaktadır.
\end{vurgukutu}

\section{Şartname Uyum Tablosu}

\begin{longtable}{p{5.2cm} c p{7.4cm}}
\toprule
\textbf{Şartname Maddesi} & \textbf{Durum} & \textbf{Eksik} \\
\midrule
\endhead
5.1 Veri Toplama & \textcolor{kritik}{Yok} & BDDK listesindeki bankaların sitelerinden toplanmış veri seti yok; elde 6 adet elle yazılmış örnek var \\
5.2 Metin Analizi Yeteneği & \textcolor{tamam}{İyi} & Dolaylı ifadeler (``avantajlı kâr payı fırsatı'') sistematik test edilmedi \\
5.3 Finansal Bilgi Çıkarımı & \textcolor{onemli}{Kısmi} & Banka adı, avantajlar, kampanya koşulları, indirim oranı, alışveriş puanı, masraf bilgisi alanları yok \\
5.4 Kampanya Türü Belirleme & \textcolor{onemli}{Kısmi} & Yalnızca \texttt{urun\_turu} var; ayrı \texttt{kampanya\_turu} sınıflandırması yok \\
5.5 Terminoloji Uyumu & \textcolor{tamam}{Güçlü} & --- \\
5.6 Standart Formata Dönüştürme & \textcolor{tamam}{Güçlü} & --- \\
5.7 Ürünlerin Karşılaştırılması & \textcolor{onemli}{Kısmi} & ``En düşük kâr payı / en yüksek ödül / en avantajlı'' sıralama mantığı yok \\
5.8 Veri Ön İşleme & \textcolor{kritik}{Yok} & Ham metin doğrudan modele gidiyor; dokümante edilebilir ön işleme hattı yok \\
5.9 On-Premise & \textcolor{kritik}{Yok} & Bulut API bağımlılığı \\
5.10 Açık Kaynak & \textcolor{onemli}{Riskli} & Kod açık fakat harici servise bağımlı; LICENSE dosyası yok \\
Dashboard & \textcolor{onemli}{Kısmi} & İnceleme paneli var, raporlama dashboard'u yok \\
Chatbot & \textcolor{kritik}{Yok} & Hiç başlanmadı \\
\bottomrule
\end{longtable}

\section{Eksiklik Listesi}

\subsection{Kritik Eksikler}

\begin{enumerate}[label=\textbf{K\arabic*.},leftmargin=1.4cm,itemsep=4pt]
  \item \textbf{Lokal çalışma yok.} Model çağrısı buluta gidiyor, banka metni kurum
        dışına çıkıyor. \hfill \kritikrozet
  \item \textbf{Chatbot yok.} Şartnamede iki örnek diyalog tanımlı: tek banka sorgusu
        (``A Bankası'nın konut finansmanı oranı ne?'') ve iki banka karşılaştırması
        (``A Bankası mı daha avantajlı, C Bankası mı?'').
  \item \textbf{Gerçek veri seti yok.} BDDK'nın katılım bankaları listesindeki
        kuruluşların resmî sitelerinden toplanmış kampanya metni bulunmuyor.
  \item \textbf{Banka adı alanı şemada yok.} Banka bilgisi çıkarılmadığı için
        bankalar arası karşılaştırma teknik olarak eksik kalıyor.
  \item \textbf{Kalıcı veri deposu yok.} Tüm veri bellekte; sayfa yenilenince
        kayboluyor. Dashboard ve chatbot'un sorgulayacağı bir kaynak yok.
  \item \textbf{Veri ön işleme adımı yok.} Normalizasyon, gürültü temizleme,
        cümle bölütleme gibi adımlar ayrı bir katman olarak mevcut değil.
\end{enumerate}

\subsection{Önemli Kapsam Eksikleri}

\begin{enumerate}[label=\textbf{Ö\arabic*.},leftmargin=1.4cm,itemsep=4pt]
  \item \textbf{Eksik çıkarım alanları:} avantajlar listesi, kampanya koşulları,
        indirim oranı, alışveriş puanı, genel masraf bilgisi. \hfill \onemlirozet
  \item \textbf{Kampanya türü sınıflandırması yok.} Şartname sekiz kampanya türü
        saymaktadır; ``Yeni Müşteri Kampanyası'' gibi türler karşılanmıyor.
  \item \textbf{Karşılaştırma kriterleri uygulanmıyor.} En düşük kâr payı, en yüksek
        ödül, en uzun vade, en düşük masraf, en avantajlı kampanya.
  \item \textbf{Dashboard yok.} Banka ve ürün bazlı toplu raporlama, dağılım
        grafikleri, oran karşılaştırması sunulmuyor.
  \item \textbf{Dolaylı ifade testi yapılmadı.} \%30 ağırlıklı ``Model Başarısı ve
        Anlamlandırma Yeteneği'' kriterinin merkezinde bu yetenek yer alıyor.
  \item \textbf{Performans ölçümü yok.} Altın veri seti ve doğruluk metriği olmadan
        model başarısı savunulamıyor.
  \item \textbf{Örnek metinler kurgu.} Gerçek banka adlarına atfedilmiş uydurma
        metinler, teslim edilecek veri seti olarak sorunludur.
\end{enumerate}

\subsection{Teslim Gereklilikleri}

\begin{enumerate}[label=\textbf{T\arabic*.},leftmargin=1.4cm,itemsep=4pt]
  \item \textbf{README yok.} Kurulum adımları, eksiksiz bağımlılık listesi ve veri
        setinin indirilebileceği herkese açık bağlantı zorunludur. \hfill \teslimrozet
  \item \textbf{LICENSE yok.} Apache License 2.0 zorunludur.
  \item \textbf{GitHub yüklemesi yapılmadı.} \texttt{BilisimVadisi2026} etiketi ve
        en az haftalık güncelleme zorunludur.
  \item \textbf{Proje dokümantasyonu yok.} Mimari, veri akışı, NLP yaklaşımı,
        ön işleme adımları, karşılaştırma yöntemi, karşılaşılan problemler,
        çıktı örnekleri, performans değerlendirme yöntemi.
  \item \textbf{Demo videosu yok.} Maksimum 5 dakika; arayüz, dashboard, chatbot,
        metin girişi, yapılandırılmış çıktı ve karşılaştırma gösterilmelidir.
        Sunum için ayrıca 1 dakikalık demo videosu gerekmektedir.
  \item \textbf{Sunum materyali yok.} PDF ve PPTX formatında, 4 dakikalık sunum.
\end{enumerate}

\subsection{Mühendislik Hijyeni}

\begin{enumerate}[label=\textbf{M\arabic*.},leftmargin=1.4cm,itemsep=4pt]
  \item Otomatik test yok (ne kural tabanlı çıkarıcı ne API sözleşmesi için).
  \item Çalışma zamanı şema doğrulaması yok; model bozuk JSON döndüğünde tipler
        gerçeği yansıtmıyor.
  \item Model çağrısında zaman aşımı ve yeniden deneme mekanizması yok.
  \item Kural tabanlı yedek çıkarıcıda olası hata: ``aylık kâr payı'' ifadesi geçen
        metinde \texttt{periyot} alanı ``belirsiz'' dönmektedir.
  \item Gerçek bir linter yapılandırması yok; \texttt{lint} betiği yalnızca tip
        denetimi yapıyor.
\end{enumerate}

\section{Puan Etkisi Tahmini}

\begin{table}[h]
\centering
\begin{tabular}{p{6.4cm} c c c}
\toprule
\textbf{Değerlendirme Kriteri} & \textbf{Ağırlık} & \textbf{Mevcut} & \textbf{Hedef} \\
\midrule
Fonksiyonellik ve Senaryo Kapsamı & \%20 & $\sim$8 & 18 \\
Teknik İmplementasyon ve Mimari & \%20 & $\sim$12 & 18 \\
On-Prem Uygulanabilirlik & \%20 & $\sim$2 & 18 \\
Model Başarısı ve Anlamlandırma & \%30 & $\sim$19 & 26 \\
Yenilikçilik ve Yaratıcılık & \%10 & $\sim$7 & 9 \\
\midrule
\textbf{Toplam} & \textbf{\%100} & $\sim$\textbf{48} & \textbf{89} \\
\bottomrule
\end{tabular}
\caption{Kaba puan tahmini. Kayıpların büyük bölümü model kalitesinden değil,
yapısal eksiklerden kaynaklanmaktadır.}
\end{table}

\section{Yol Haritası}

\subsection{Faz 1 --- Temel (1--2 gün)}

\textbf{Amaç:} En yüksek puan kaybını en düşük eforla durdurmak.

\begin{enumerate}[leftmargin=1.2cm,itemsep=3pt]
  \item \textbf{Lokal LLM'e geçiş.} Ollama veya vLLM üzerinde açık kaynak bir model
        (Qwen2.5, Llama 3.1 8B vb.) çalıştırılır. Mevcut kod OpenAI uyumlu bir
        endpoint'e istek attığı için taban adresin değiştirilmesi büyük ölçüde
        yeterlidir. Bulut sağlayıcı isteğe bağlı yedek konumuna alınır.
        \emph{(Çözer: K1, T-riski)}
  \item \textbf{Kalıcı depo.} SQLite tabanlı, çıkarılan ürünleri, kaynak metni ve
        analist doğrulamalarını saklayan katman. \emph{(Çözer: K5)}
  \item \textbf{Şema genişletme.} \texttt{banka\_adi}, \texttt{kampanya\_turu},
        \texttt{avantajlar[]}, \texttt{kampanya\_kosullari}, \texttt{indirim\_orani},
        \texttt{alisveris\_puani} alanları eklenir. \emph{(Çözer: K4, Ö1, Ö2)}
\end{enumerate}

\subsection{Faz 2 --- Senaryo Kapsamı (2--3 gün)}

\begin{enumerate}[leftmargin=1.2cm,itemsep=3pt]
  \item \textbf{Chatbot.} Doğal dildeki soruyu veritabanı sorgusuna çeviren
        araç çağırma (tool calling) yaklaşımı. Modelin serbest üretimi yerine
        yapılandırılmış veriden cevap üretilir; cevabın altına kanıt alıntısı
        eklenir. Şartnamedeki iki diyalog senaryosu birebir karşılanır.
        \emph{(Çözer: K2)}
  \item \textbf{Karşılaştırma motoru.} En düşük kâr payı, en yüksek ödül, en uzun
        vade, en düşük masraf ve bileşik ``en avantajlı'' sıralaması.
        \emph{(Çözer: Ö3)}
  \item \textbf{Dashboard.} Banka ve ürün türü bazlı dağılım, kâr payı oranı
        karşılaştırması, manuel inceleme bekleyen kayıt sayacı, son çıkarımlar.
        \emph{(Çözer: Ö4)}
\end{enumerate}

\subsection{Faz 3 --- Veri ve Doğruluk (2--3 gün)}

\begin{enumerate}[leftmargin=1.2cm,itemsep=3pt]
  \item \textbf{Veri toplama hattı.} BDDK listesindeki katılım bankalarının kampanya
        sayfaları için Python tabanlı toplayıcı. Tam otomatik tarama yerine banka
        başına tanımlı kaynak URL listesi daha az kırılgan ve savunması daha
        kolaydır. Toplanan veri seti açık bağlantı ile paylaşılır.
        \emph{(Çözer: K3, Ö7, T1)}
  \item \textbf{Ön işleme hattı.} HTML temizleme, boşluk ve tırnak normalizasyonu,
        cümle bölütleme, para birimi ve sayı biçimi standardizasyonu. Ayrı modül
        olarak yazılır ki dokümantasyonda gösterilebilsin. \emph{(Çözer: K6)}
  \item \textbf{Değerlendirme seti.} Elle doğrulanmış altın JSON'lar ve alan bazlı
        kesinlik/duyarlılık ölçen betik. Dolaylı ifade örnekleri özellikle bu sete
        dâhil edilir. \emph{(Çözer: Ö5, Ö6)}
\end{enumerate}

\subsection{Faz 4 --- Teslim (1--2 gün)}

\begin{enumerate}[leftmargin=1.2cm,itemsep=3pt]
  \item README, LICENSE (Apache 2.0), proje dokümantasyonu.
  \item GitHub deposu, \texttt{BilisimVadisi2026} etiketi, haftalık güncelleme
        düzeni.
  \item Demo videosu (5 dk) ve sunum videosu (1 dk).
  \item Sunum materyali (PDF + PPTX, 4 dk).
  \item Test, şema doğrulama, zaman aşımı ve linter düzenlemeleri.
        \emph{(Çözer: T1--T6, M1--M5)}
\end{enumerate}

\subsection{Bağımlılık Sırası}

Fazlar arasında zorunlu sıra bulunmaktadır:

\begin{center}
\texttt{Kalıcı depo} $\rightarrow$ \texttt{Chatbot + Dashboard}\\[2pt]
\texttt{Banka adı alanı} $\rightarrow$ \texttt{Karşılaştırma motoru}\\[2pt]
\texttt{Veri toplama} $\rightarrow$ \texttt{Değerlendirme seti}
\end{center}

\section{Korunması Gereken Güçlü Yanlar}

Mevcut çözümün üç ayırt edici özelliği vardır ve geliştirme sürecinde
kaybedilmemelidir:

\begin{itemize}[leftmargin=1.2cm,itemsep=3pt]
  \item \textbf{Kanıt alıntısı.} Her alan için, değerin dayandığı cümle kaynak
        metinden birebir alıntılanmaktadır. Chatbot cevaplarına da eklenmelidir.
  \item \textbf{Alan bazlı güven skoru.} Düşük güvenli alanlar işaretlenip insan
        doğrulamasına yönlendirilmektedir.
  \item \textbf{İnsan doğrulama döngüsü.} Analistin alan bazında onay verebildiği
        arayüz, bankacılık gibi denetime tabi bir alanda gerçek hayata
        uygulanabilirliği güçlendirmektedir.
\end{itemize}

\section{Doğrulanması Gereken Husus}

Şartname içinde takvim tutarsızlığı bulunmaktadır: Tablo 1'de yarışma çevrimiçi
süreci 27 Temmuz --- 26 Ağustos olarak verilirken, Madde 9'da yarışmanın 12.07.2026
tarihinde sona ereceği ve proje dosyalarının bu tarihe kadar yüklenmesi gerektiği
belirtilmektedir. Planlama yapılmadan önce geçerli tarihin TEKNOFEST duyuru kanalları
üzerinden teyit edilmesi gerekmektedir.

\end{document}
